\paragraph{escort Rényi 熵常数的 \texorpdfstring{$\theta(q)$}{theta(q)}-闭式化}
在 Rényi 熵率塌缩与 escort 常数谱已经确立之后，温度参数 \(q\) 的全部 \(O(1)\) 熵信息还可继续压缩为末位 Bernoulli 质量
\(
\theta(q)=1/(1+\varphi^{q+1})
\)。

\begin{theorem}[escort Rényi 熵的速率塌缩与常数项全解]\label{thm:xi-fold-escort-renyi-constant-theta-closed}
\leanverified{paper\_xi\_fold\_escort\_renyi\_constant\_theta\_closed}
沿用定理 \ref{thm:xi-fold-escort-renyi-shannon-constants} 的记号。记
\[
\theta(q):=\frac{1}{1+\varphi^{q+1}}.
\]
则对任意固定 \(q\ge 0\) 与任意 \(r>0\)、\(r\neq 1\)，有
\[
H_r(\mu_{m,q})
=
m\log\varphi
-\frac12\log 5
+
\frac{1}{1-r}
\log\!\Bigl(
\varphi^{1-r}(1-\theta(q))^r+\theta(q)^r
\Bigr)
+
O\!\Bigl(\Bigl(\frac{\varphi}{2}\Bigr)^m\Bigr).
\]
等价地，
\[
H_r(\mu_{m,q})
=
m\log\varphi
-\frac12\log 5
+
\frac{1}{1-r}
\Bigl[
\log\!\bigl(1+\varphi^{1+qr}\bigr)
-r\log\!\bigl(1+\varphi^{q+1}\bigr)
\Bigr]
+
O\!\Bigl(\Bigl(\frac{\varphi}{2}\Bigr)^m\Bigr).
\]
相应的 Shannon 熵满足
\[
H(\mu_{m,q})
=
m\log\varphi
-\frac12\log 5
+
h_2(\theta(q))
+
(1-\theta(q))\log\varphi
+
O\!\Bigl(\Bigl(\frac{\varphi}{2}\Bigr)^m\Bigr),
\]
其中
\[
h_2(p):=-p\log p-(1-p)\log(1-p).
\]
因此 escort 温度 \(q\) 的全部 \(O(1)\) 熵信息都被压缩为单变量 \(\theta(q)\)。
\end{theorem}
\begin{proof}
由定理 \ref{thm:xi-fold-escort-renyi-shannon-constants}，
\[
H_r(\mu_{m,q})
=
m\log\varphi
+
\frac{1}{1-r}\left[
\log\!\Bigl(\frac{\varphi+\varphi^{-qr}}{\sqrt5}\Bigr)
-r\log\!\Bigl(\frac{\varphi+\varphi^{-q}}{\sqrt5}\Bigr)
\right]
+
O\!\Bigl(\Bigl(\frac{\varphi}{2}\Bigr)^m\Bigr).
\]
再用
\[
\varphi+\varphi^{-a}=\varphi^{-a}(1+\varphi^{a+1}),
\]
便得到第二个 Rényi 公式。

另一方面，由
\[
1-\theta(q)=\frac{\varphi^{q+1}}{1+\varphi^{q+1}}
\]
可直接化简出
\[
\varphi^{1-r}(1-\theta(q))^r+\theta(q)^r
=
\frac{1+\varphi^{1+qr}}{(1+\varphi^{q+1})^r},
\]
故第一个 Rényi 公式与第二个完全等价。

对 Shannon 熵，定理 \ref{thm:xi-fold-escort-renyi-shannon-constants} 同样给出
\[
H(\mu_{m,q})
=
m\log\varphi
+
\log\!\Bigl(\frac{\varphi+\varphi^{-q}}{\sqrt5}\Bigr)
+
\frac{q\log\varphi}{1+\varphi^{q+1}}
+
O\!\Bigl(\Bigl(\frac{\varphi}{2}\Bigr)^m\Bigr).
\]
再注意到
\[
h_2(\theta(q))
=
\log\!\bigl(1+\varphi^{q+1}\bigr)
-(q+1)(1-\theta(q))\log\varphi,
\]
于是
\[
-\frac12\log 5
+
h_2(\theta(q))
+
(1-\theta(q))\log\varphi
=
\log\!\Bigl(\frac{\varphi+\varphi^{-q}}{\sqrt5}\Bigr)
+
\frac{q\log\varphi}{1+\varphi^{q+1}}.
\]
这就给出了 Shannon 常数项的所述写法。证毕。
\end{proof}

\endinput
